\documentclass{beamer}

\usepackage[T2A]{fontenc}
\usepackage[utf8]{inputenc}
\usepackage[english,russian]{babel}
\pdfmapfile{+sansmathaccent.map}


\mode<presentation>
{
	\usetheme{Warsaw} % or try Darmstadt, Madrid, Warsaw, Rochester, CambridgeUS, ...
	\usecolortheme{seahorse} % or try seahorse, beaver, crane, wolverine, ...
	\usefonttheme{serif}  % or try serif, structurebold, ...
	\setbeamertemplate{navigation symbols}{}
	\setbeamertemplate{caption}[numbered]
} 


%%%%%%%%%%%%%%%%%%%%%%%%%%%%
% itemize settings


%%%%%%%%%%%%%%%%%%%%%%%%%%%%
% itemize settings

\definecolor{myhotpink}{RGB}{255, 80, 200}
\definecolor{mywarmpink}{RGB}{255, 60, 160}
\definecolor{mylightpink}{RGB}{255, 80, 200}
\definecolor{mypink}{RGB}{255, 30, 80}
\definecolor{mydarkpink}{RGB}{155, 25, 60}

\definecolor{mypaleblue}{RGB}{240, 240, 255}
\definecolor{mylightblue}{RGB}{120, 150, 255}
\definecolor{myblue}{RGB}{90, 90, 255}
\definecolor{mygblue}{RGB}{70, 110, 240}
\definecolor{mycourseblue}{RGB}{120, 140, 255}
\definecolor{mycourselightblue}{RGB}{245, 250, 255}
\definecolor{mydarkblue}{RGB}{0, 0, 180}
\definecolor{myblackblue}{RGB}{40, 40, 120}

\definecolor{mygreen}{RGB}{0, 200, 0}
\definecolor{mygreen2}{RGB}{245, 255, 230}
\definecolor{mydarkgreen}{RGB}{0, 120, 0}
\definecolor{mydarkgreen}{RGB}{0, 120, 0}


\definecolor{mygray}{gray}{0.8}
\definecolor{mydarkgray}{RGB}{80, 80, 160}
\definecolor{mylightgray}{RGB}{160, 160, 160}

\definecolor{mydarkred}{RGB}{160, 30, 30}
\definecolor{mylightred}{RGB}{255, 150, 150}
\definecolor{myred}{RGB}{200, 110, 110}
\definecolor{myblackred}{RGB}{120, 40, 40}

\definecolor{mypink}{RGB}{255, 30, 80}
\definecolor{myhotpink}{RGB}{255, 80, 200}
\definecolor{mywarmpink}{RGB}{255, 60, 160}
\definecolor{mylightpink}{RGB}{255, 80, 200}
\definecolor{mydarkpink}{RGB}{155, 25, 60}
\definecolor{mywhitepink}{RGB}{255, 240, 240}

\definecolor{mydarkcolor}{RGB}{60, 25, 155}
\definecolor{mylightcolor}{RGB}{130, 180, 250}



\setbeamertemplate{itemize items}[default]

\setbeamertemplate{itemize item}{\color{myblackblue}$\blacksquare$}
\setbeamertemplate{itemize subitem}{\color{mygblue}$\blacktriangleright$}
\setbeamertemplate{itemize subsubitem}{\color{mygray}$\blacksquare$}

\setbeamercolor{palette quaternary}{fg=white,bg=mycourseblue}
\setbeamercolor{titlelike}{parent=palette quaternary}

\setbeamercolor{palette quaternary2}{fg=black,bg=mycourselightblue}
\setbeamercolor{frametitle}{parent=palette quaternary2}

\setbeamerfont{frametitle}{size=\Large,series=\scshape}
\setbeamerfont{framesubtitle}{size=\normalsize,series=\upshape}





%%%%%%%%%%%%%%%%%%%%%%%%%%%%
% block settings

\setbeamercolor{block title}{bg=red!30,fg=black}

\setbeamercolor*{block title example}{bg=mygreen!40!white,fg=black}

\setbeamercolor*{block body example}{fg= black, bg= mygreen2}


%%%%%%%%%%%%%%%%%%%%%%%%%%%%
% URL settings
\hypersetup{
	colorlinks=true,
	linkcolor=blue,
	filecolor=blue,      
	urlcolor=blue,
}

%%%%%%%%%%%%%%%%%%%%%%%%%%

\renewcommand{\familydefault}{\rmdefault}


\newcommand{\mycolor}[1] {
\setbeamercolor{normal text}{fg={#1}} % Recolors standard body text
\setbeamercolor{structure}{fg={#1}}   % Recolors bullets and structure items
\usebeamercolor[fg]{normal text}      % Applies the normal text color immediately
}


\usepackage{amsmath}
\usepackage{mathtools}

\usepackage{cancel}
\usepackage{centernot}

\usepackage{changepage}
\usepackage{subcaption}

\usepackage{qrcode}

\newcommand{\bo}[1] {\mathbf{#1}}
\newcommand{\R}{\mathbb{R}} 
\newcommand{\T}{^\top}     



\newcommand{\mydate}{Fall 2026}

\newcommand{\mygit}{\textcolor{blue}{\href{https://github.com/SergeiSa/Advanced-Control-2026}{github.com/SergeiSa/Advanced-Control-2026}}}

\newcommand{\myqr}{ \textcolor{black}{\qrcode[height=1.5in]{https://github.com/SergeiSa/Advanced-Control-2026}}
}




\newcommand{\bref}[2] {\textcolor{blue}{\href{#1}{#2}}}

%%%%%%%%%%%%%%%%%%%%%%%%%%%%
% code settings

\usepackage{listings}
\usepackage{color}
% \definecolor{mygreen}{rgb}{0,0.6,0}
% \definecolor{mygray}{rgb}{0.5,0.5,0.5}
\definecolor{mymauve}{rgb}{0.58,0,0.82}
\lstset{ 
	backgroundcolor=\color{white},   % choose the background color; you must add \usepackage{color} or \usepackage{xcolor}; should come as last argument
	basicstyle=\footnotesize,        % the size of the fonts that are used for the code
	breakatwhitespace=false,         % sets if automatic breaks should only happen at whitespace
	breaklines=true,                 % sets automatic line breaking
	captionpos=b,                    % sets the caption-position to bottom
	commentstyle=\color{mygreen},    % comment style
	deletekeywords={...},            % if you want to delete keywords from the given language
	escapeinside={\%*}{*)},          % if you want to add LaTeX within your code
	extendedchars=true,              % lets you use non-ASCII characters; for 8-bits encodings only, does not work with UTF-8
	firstnumber=0000,                % start line enumeration with line 0000
	frame=single,	                   % adds a frame around the code
	keepspaces=true,                 % keeps spaces in text, useful for keeping indentation of code (possibly needs columns=flexible)
	keywordstyle=\color{blue},       % keyword style
	language=Octave,                 % the language of the code
	morekeywords={*,...},            % if you want to add more keywords to the set
	numbers=left,                    % where to put the line-numbers; possible values are (none, left, right)
	numbersep=5pt,                   % how far the line-numbers are from the code
	numberstyle=\tiny\color{mygray}, % the style that is used for the line-numbers
	rulecolor=\color{black},         % if not set, the frame-color may be changed on line-breaks within not-black text (e.g. comments (green here))
	showspaces=false,                % show spaces everywhere adding particular underscores; it overrides 'showstringspaces'
	showstringspaces=false,          % underline spaces within strings only
	showtabs=false,                  % show tabs within strings adding particular underscores
	stepnumber=2,                    % the step between two line-numbers. If it's 1, each line will be numbered
	stringstyle=\color{mymauve},     % string literal style
	tabsize=2,	                   % sets default tabsize to 2 spaces
	title=\lstname                   % show the filename of files included with \lstinputlisting; also try caption instead of title
}


%%%%%%%%%%%%%%%%%%%%%%%%%%%%
% URL settings
\hypersetup{
	colorlinks=false,
	linkcolor=blue,
	filecolor=blue,      
	urlcolor=blue,
}

%%%%%%%%%%%%%%%%%%%%%%%%%%

%%%%%%%%%%%%%%%%%%%%%%%%%%%%
% tikz settings

\usepackage{tikz}
\tikzset{every picture/.style={line width=0.75pt}}
\usetikzlibrary{shapes.geometric, decorations.markings, positioning, patterns, calc, arrows.meta}

\newcommand{\myqrframe}{
	\begin{frame}
		\begin{center}
			\centerline{Слайды доступны на Github:}
			\bigskip
			\centerline{\mygit}
			\bigskip
			\myqr
		\end{center}
	\end{frame}
}


\subtitle{Современные методы управления роботами}
\author{Сергей Савин}
\date{2026}

\title{Планирование траектории}



\begin{document}
	\maketitle
	
	
	\begin{frame}{Содержание}
		\begin{itemize}
			\item Устойчивость и планирование траектории
			\item Планирование траектории — иллюстрация
			\item Пример: раскачивание перевёрнутого маятника
			\item Пример: тележка с шестом
			\item Путь vs.\ траектория
			\item Допустимая траектория
			\item Оптимизация траектории
			\item Прямой пристрелочный метод
			\item Параметризация управляющего воздействия
			\item Решение задачи прямого пристрелочного метода
			\item Литература
		\end{itemize}
	\end{frame}
	
	
	\begin{frame}{Устойчивость и планирование траектории}
		\begin{flushleft}
			
			Для динамической системы
			\[
			\dot{x}=f(x,u),
			\]
			можно рассмотреть две различные задачи:
			
			\begin{enumerate}
				\item Найти закон управления, который переводит систему
				из начального состояния $x_A$ в желаемое состояние $x_B$.
				
				\item Найти закон управления, который заставляет систему
				оставаться вблизи $x_B$ (или сходиться к нему), как только она
				достаточно близко к нему приблизилась.
			\end{enumerate}
			
			\vspace{0.5em}
			
			Вторая задача — это задача
			\textbf{стабилизирующего управления}.
			
			Первая задача — это задача
			\textbf{планирования траектории}.
			
		\end{flushleft}
	\end{frame}
	
	
	\begin{frame}{Планирование траектории — иллюстрация}
		\begin{flushleft}
			
				\begin{tikzpicture}[
		% Arrow decoration for the trajectory path
		path arrow/.style={
			postaction={
				decorate,
				decoration={
					markings,
					mark=at position 0.55 with {\arrow[scale=1.3]{stealth}}
				}
			}
		},
		% Point marker style
		point/.style={
			circle,
			fill=black,
			inner sep=1.5pt
		}
		]
		
		% Define coordinates
		\coordinate (A) at (-4, -0.5);
		\coordinate (B) at (3, 0.5);
		
		% --- Region of Attraction (ROA) Fill ---
		% Fills the outer ellipsoid around point B
		\fill[blue!10, rotate around={-25:(B)}] (B) ellipse (2.5cm and 1.8cm);
		
		% --- Lyapunov Function Ellipsoids (Contours) ---
		% Diminishing sizes around point B
		\begin{scope}[rotate around={-25:(B)}]
			\draw[thin, color=gray!80] (B) ellipse (2.5cm and 1.8cm);
			\draw[thin, color=gray!80] (B) ellipse (1.9cm and 1.37cm);
			\draw[thin, color=gray!80] (B) ellipse (1.3cm and 0.94cm);
			\draw[thin, color=gray!80] (B) ellipse (0.7cm and 0.5cm);
		\end{scope}
		
		% --- Trajectory Path ---
		% Smooth curved path from A into B
		\draw[thick, dashed, path arrow] 
		(A) .. controls (-2, 2) and (0, -1.8) .. (B);
		
		% --- Points and Labels ---
		\node[point, label=left:{$A$}] at (A) {};
		\node[point, label=right:{$B$ (Fixed Point)}] at (B) {};
		
		% --- Optional Labels ---
		\node[blue!60!black, font=\small] at ($(B) + (0, 2.1)$) {Region of Attraction (ROA)};
		\node[gray!90!black, font=\footnotesize] at ($(B) + (1.6, -1.3)$) {$V(x) = c$};
		
	\end{tikzpicture}
			
		\end{flushleft}
	\end{frame}
	
	
	\begin{frame}{Пример: раскачивание перевёрнутого маятника}
		\begin{flushleft}
			
			
\begin{center}
\begin{tikzpicture}[
	scale=0.85,
	>=Stealth,
	base/.style={thick, line cap=round},
	stable/.style={draw=black, line width=1.1pt},
	intermediate/.style={draw=black!45, line width=0.9pt, dashed},
	pivot/.style={fill=white, draw=black, thick},
	bob/.style={circle, fill=black, inner sep=2.2pt},
	bob gray/.style={circle, fill=black!45, inner sep=2pt},
	arc arrow/.style={->, draw=black!70, thick, >=Stealth}
	]
	
	\def\L{2.2} % Length of pendulum arm
	
	% --- LEFT SIDE: LOWER EQUILIBRIUM ---
	\coordinate (O1) at (0, 0);
	
	% Base Mount
	\draw[base] (-0.8,0) -- (0.8,0);
	\fill[pattern=north east lines, pattern color=black!40] (-0.8,-0.15) rectangle (0.8,0);
	
	% Ground / Reference dashed line
	\draw[gray!40, dash pattern=on 2pt off 2pt] (O1) -- ++(0, -\L - 0.4);
	
	% 4 States
	\draw[stable] (O1) -- ++(270:\L) node[bob] {};
	\draw[stable] (O1) -- ++(255:\L) node[bob] {};
	\draw[intermediate] (O1) -- ++(225:\L) node[bob gray] {};
	\draw[intermediate] (O1) -- ++(205:\L) node[bob gray] {};
	
	% Pivot Joint
	\draw[pivot] (O1) circle (3pt);
	
	% Arc Arrow
	\draw[arc arrow] (265:1.4) arc (265:210:1.4);
	
	
	% --- RIGHT SIDE: UPPER EQUILIBRIUM (Shifted down to align top-to-top) ---
	% Offset y by -\L so the top of the 180-deg pendulum reaches y = 0
	\coordinate (O2) at (4.5, -\L);
	
	% Base Mount
	\draw[base] (O2) ++(-0.8,0) -- ++(1.6,0);
	\fill[pattern=north east lines, pattern color=black!40] (O2) ++(-0.8,-0.15) rectangle ++(1.6,0.15);
	
	% Vertical Reference line
	\draw[gray!40, dash pattern=on 2pt off 2pt] (O2) -- ++(0, \L + 0.4);
	
	% States
	\draw[intermediate] (O2) -- ++(135:\L) node[bob gray] {};
	\draw[stable] (O2) -- ++(170:\L) node[bob] {};
	\draw[stable] (O2) -- ++(90:\L) node[bob] {};
	
	% Pivot Joint
	\draw[pivot] (O2) circle (3pt);
	
	% Arc Arrow
	\draw[arc arrow] ($(O2)+(138:1.5)$) arc (138:95:1.5);
	
\end{tikzpicture}
\end{center}
			
			
			Для маятника можно рассмотреть две отдельные задачи: 1) спланировать траекторию, которая переводит маятник из нижнего положения равновесия в верхнее; 2) найти закон управления, который стабилизирует верхнее положение равновесия.
			
			\bigskip
			
			Первая задача может быть сложной, если двигатель слабый
			и имеет малый предел крутящего момента. Вместо того чтобы поднимать маятник напрямую, возможно, потребуется
			\textbf{раскачать его}.
			
		\end{flushleft}
	\end{frame}
	
	
	
	
	\begin{frame}{Пример: тележка с шестом}
		\begin{flushleft}
			
			% TODO: \usepackage{graphicx} required
			\begin{figure}
				\centering
				\includegraphics[width=0.7\linewidth]{"cartpole swing up"}
				\label{fig:cartpole-swing-up}
			\end{figure}
			
			
			Для тележки с шестом мы снова можем рассмотреть две отдельные задачи: 1) спланировать траекторию, которая переводит шест
			из нижнего положения равновесия в верхнее; 2) найти закон управления, который стабилизирует верхнее положение равновесия.
			
			\medskip
			
			Между тележкой и шестом нет двигателя напрямую. Поэтому мы не можем просто приложить крутящий момент: вместо этого мы должны двигать тележку так, чтобы её движение создавало \textbf{инерционные силы}, которые раскачивают шест вверх.
			
		\end{flushleft}
	\end{frame}
	
	
	\begin{frame}{Путь vs.\ траектория}
		\begin{flushleft}
			
			Рассмотрим два состояния $x_A$ и $x_B$. Кривая, соединяющая их в пространстве состояний, называется
			\textbf{путём}.
			
			\bigskip
			
			Путь говорит нам о последовательности состояний, через которые система должна пройти, но не указывает, \emph{когда} эти состояния должны быть достигнуты.
			
			\bigskip
			
			Путь вместе с его временной развёрткой называется	\textbf{траекторией}. Мы можем обозначить её как $x(t)$.
			
		\end{flushleft}
	\end{frame}
	
	
	\begin{frame}{Допустимая траектория}
		\begin{flushleft}
			
			Не каждая траектория может быть фактически получена динамической системой. Для системы с управляющим воздействием
			\[
			\dot{x}=f(x,u),
			\]
			траектория $x^*(t)$ является \emph{допустимой}, если существует управляющее воздействие $u^*(t)$ такое, что
			
			\[
			\dot{x}^*(t)
			=
			f\bigl(x^*(t),u^*(t)\bigr),
			\qquad
			\forall t\in[0,T].
			\]
			
			
			\bigskip
			
			\begin{center}
				\fbox{Траектория должна удовлетворять динамике системы.}
			\end{center}
			
		\end{flushleft}
	\end{frame}
	
	
	
	\begin{frame}{Оптимизация траектории}
		\begin{flushleft}
			Оптимизация траектории — это задача поиска
			\emph{допустимой траектории}, которая минимизирует выбранную
			функцию стоимости.
			
			\vspace{0.5em}
			
			Например, мы можем сформулировать задачу как
			\[
			\begin{aligned}
				\min_{x(t),u(t),T}\quad
				&
				\Phi\bigl(x(T)\bigr)
				+
				\int_0^T L\bigl(x(t),u(t)\bigr)\,dt
				\\
				\text{s.t.}\quad
				&
				\dot{x}(t)=f\bigl(x(t),u(t)\bigr),
				\\
				&
				x(0)=x_A,
				\qquad
				x(T)=x_B.
			\end{aligned}
			\]
			
			\begin{itemize}
				\item $\Phi(x(T))$ — это \textbf{терминальная стоимость}, которая
				штрафует конечное состояние.
				
				\item $L(x(t),u(t))$ — это \textbf{стоимость процесса}, которая
				штрафует состояние и управление вдоль траектории.
			\end{itemize}
			
			Конкретный выбор $\Phi$ и $L$ зависит от того, чего
			мы хотим достичь с помощью траектории.
		\end{flushleft}
	\end{frame}
	
	\begin{frame}{Прямой пристрелочный метод}
		\begin{flushleft}
			Заметим, что для заданного управляющего воздействия $u(t)$ и
			начального условия $x(0)=x_0$ результирующая траектория полностью определяется
			динамикой системы
			\[
			\dot{x}=f(x,u).
			\]
			
			\vspace{0.5em}
			
			Следовательно, нам не нужно выбирать и $u(t)$, и $x(t)$.
			
			Вместо этого мы можем:
			\begin{itemize}
				\item выбрать $u(t)$;
				\item просимулировать систему вперёд, чтобы получить $x(t)$.
			\end{itemize}
			
			\vspace{0.5em}
			
			Этот подход к оптимизации траектории называется
			\textbf{прямым пристрелочным методом}.
		\end{flushleft}
	\end{frame}
	
	
	\begin{frame}{Параметризация управляющего воздействия}
		\begin{flushleft}
			Существует несколько способов параметризовать $u(t)$ для
			прямого пристрелочного метода.
			
			\begin{itemize}
				\item Кусочно-постоянная:
				\[
				u(t)=u_i,
				\qquad
				t\in[t_i,t_{i+1}).
				\]
				
				\item Полиномы более высокого порядка на каждом временном интервале.
				
				\item Сплайны.
				
				\item Единая функция, зависящая от времени, например:
				\[
				u(t)=c_0+c_1t+c_2t^2+c_3t^3.
				\]
			\end{itemize}
			
			\vspace{0.5em}
			
			Ключевая идея состоит в том, чтобы параметризовать управляющее воздействие с помощью
			конечного числа параметров $c$:
			\[
			u(t)=u(c,t).
			\]
			
			Параметры $c$ затем становятся
			\textbf{переменными решения} задачи оптимизации.
		\end{flushleft}
	\end{frame}
	
	
	\begin{frame}{Решение задачи прямого пристрелочного метода}
		\begin{flushleft}
			После параметризации управляющего воздействия как
			\[
			u(t)=u(c,t),
			\]
			задача оптимизации траектории становится задачей оптимизации
			по конечномерному вектору $c$.
			
			\vspace{0.5em}
			
			Поясняющая (хотя и неполная) картина
			процесса оптимизации такова:
			
			\begin{enumerate}
				\item По текущей оценке $c$ построить
				$u(c,t)$.
				
				\item Решить $\dot{x}=f(x,u(c,t))$ в прямом времени.
				
				\item Вычислить стоимость и её градиент $J(c),
				\quad
				\nabla_c J(c)$.
				
				\item Обновить $c$ в направлении, которое уменьшает стоимость.
				
				\item Повторять до сходимости.
			\end{enumerate}
			
			\vspace{0.5em}
\begin{adjustwidth}{-0.8cm}{-0.8cm}
	Это даёт траекторию $x(t)$ и управляющее воздействие $u(t)$,
	которые приближённо решают задачу оптимизации траектории.
\end{adjustwidth}
		\end{flushleft}
	\end{frame}

\begin{frame}{Литература}
	
	\begin{itemize}
		
		
		\item Russ Tedrake -- MIT:
		\bref{https://underactuated.mit.edu/trajopt.html}
		{Underactuated Robotics: Trajectory Optimization, ch. 10.}
		
	\end{itemize}
	
\end{frame}



\myqrframe




\end{document}
